\documentclass{ximera}

\input{../../../preamble.tex}


\definecolor{fvdnavy}{HTML}{071D43}
\definecolor{fvdcyan}{HTML}{20BFC6}
\definecolor{fvdcyanlight}{HTML}{EFFCFB}
\definecolor{fvdmagenta}{HTML}{F10D69}
\definecolor{fvdmagentalight}{HTML}{FFF1F7}
\definecolor{fvdtext}{HTML}{163044}





%=================================================
% Pedagogische omgevingen
% PDF: tcolorbox
% HTML: learning-box
%=================================================

\newenvironment{voorkennis}
{%
    \ifdefined\HCode
        \HCode{%
<section class="learning-box learning-box--prior">
<div class="learning-box__header">
<span class="learning-box__icon" aria-hidden="true"></span>
<span class="learning-box__title">Voorkennis</span>
</div>
<div class="learning-box__content">
        }%
        \begin{itemize}
    \else
       \begin{tcolorbox}[
    enhanced,
    breakable,
    colback=fvdcyanlight,
    colframe=fvdcyan,
    coltext=fvdtext,
    boxrule=0.5pt,
    leftrule=4pt,
    arc=4mm,
    outer arc=4mm,
    left=7mm,
    right=6mm,
    top=5mm,
    bottom=5mm,
    before skip=6mm,
    after skip=6mm,
    title=Voorkennis,
    coltitle=fvdcyan,
    fonttitle=\bfseries\Large,
    attach title to upper={\par\smallskip},
    boxed title style={
        empty,
        left=0mm,
        right=0mm,
        top=0mm,
        bottom=0mm
    },
    overlay={
        \draw[
            fvdcyan,
            line width=0.5pt
        ]
        ([xshift=42mm,yshift=-13mm]frame.north west)
        --
        ([xshift=-7mm,yshift=-13mm]frame.north east);
    }
]
        \begin{itemize}
    \fi
}
{%
    \end{itemize}
    \ifdefined\HCode
        \HCode{%
</div>
</section>
        }%
    \else
        \end{tcolorbox}
    \fi
}


\newenvironment{leerdoelen}
{%
    \ifdefined\HCode
        \HCode{%
<section class="learning-box learning-box--goals">
<div class="learning-box__header">
<span class="learning-box__icon" aria-hidden="true"></span>
<span class="learning-box__title">Leerdoelen</span>
</div>
<div class="learning-box__content">
        }%
        \begin{itemize}
    \else
\begin{tcolorbox}[
    enhanced,
    breakable,
    colback=fvdmagentalight,
    colframe=fvdmagenta,
    coltext=fvdtext,
    boxrule=0.5pt,
    leftrule=4pt,
    arc=4mm,
    outer arc=4mm,
    left=7mm,
    right=6mm,
    top=5mm,
    bottom=5mm,
    before skip=6mm,
    after skip=6mm,
    title=Leerdoelen,
    coltitle=fvdmagenta,
    fonttitle=\bfseries\Large,
    attach title to upper={\par\smallskip},
    boxed title style={
        empty,
        left=0mm,
        right=0mm,
        top=0mm,
        bottom=0mm
    },
    overlay={
        \draw[
            fvdmagenta,
            line width=0.5pt
        ]
        ([xshift=42mm,yshift=-13mm]frame.north west)
        --
        ([xshift=-7mm,yshift=-13mm]frame.north east);
    }
]
        \begin{itemize}
    \fi
}
{%
    \end{itemize}
    \ifdefined\HCode
        \HCode{%
</div>
</section>
        }%
    \else
        \end{tcolorbox}
    \fi
}


\addPrintStyle{../../..}

\title{Oefeningen --- Parate kennis}
\author{Ingmar Herreman}

\begin{abstract}
Deze oefeningenbundel hoort bij het thema \emph{Parate kennis --- Je wiskundige gereedschapskist}.
De bundel is bewust ruimer dan het minimumtraject. Oefeningen met ESS moet je
zeker beheersen. De overige oefeningen gebruik je om gericht te oefenen, te verdiepen of jezelf
extra uit te dagen.
\end{abstract}

\begin{document}
\fvdchapterdata{1}{Verbanden, functies en gedrag}{2}
\maketitle

\FVDsection{Hoe gebruik je deze bundel?}

Werk niet noodzakelijk van de eerste tot de laatste oefening.
Start per onderdeel met de oefeningen aangeduid met \ESS.
Lukt dat vlot en zonder hulp, dan is de basis voor dat onderdeel voldoende paraat.
Maak je fouten, kies dan extra \oefB-oefeningen. De \oefU-oefeningen vragen meer inzicht,
combinatie of interpretatie. De \oefG-oefeningen zijn uitdagender en zijn geen onderdeel van
het minimumtraject.

% ============================================================
\FVDsection{Rekenen met getallen en breuken}

\begin{oefening}[Rekenvolgorde]{\oefB\ESS}
Bereken exact.
\begin{enumerate}
\item $18-3(2+4)$
\item $5+2^3\cdot3-7$
\item $4(7-3)-2(5+1)$
\item $\dfrac{24}{3\cdot2}+5$
\item $3-[2-4(1-3)]$
\item $\dfrac{2^4-4}{3}+1$
\end{enumerate}
\end{oefening}

\begin{oefening}[Breuken]{\oefB\ESS}
Werk uit en vereenvoudig.
\begin{enumerate}
\item $\dfrac34+\dfrac56$
\item $\dfrac78-\dfrac13$
\item $\dfrac5{12}+\dfrac7{18}$
\item $\dfrac23\cdot\dfrac9{10}$
\item $\dfrac56:\dfrac{10}{9}$
\item $\dfrac{3x}{5}+\dfrac{2x}{3}$
\item $\dfrac{4x}{9}\cdot\dfrac{3}{2x^2}$
\item $\dfrac{x^2-4}{x^2+2x}$
\item $\dfrac{x^2-9}{x^2-6x+9}$
\end{enumerate}
Vermeld bij algebraïsche breuken ook de uitgesloten waarden.
\end{oefening}

\begin{oefening}[Gemengde breuken]{\oefB}
Werk zo ver mogelijk uit.
\begin{enumerate}
\item $\dfrac1x+\dfrac1{x+1}$
\item $\dfrac{x}{x+2}+\dfrac3{x+2}$
\item $\dfrac2{x-1}-\dfrac5{x+1}$
\item $\dfrac{x^2+5x+6}{x+2}$
\item $\dfrac{2x^2-8}{x+2}$
\item $\dfrac{\frac{12}{5}}{3}$
\item $\dfrac{\frac{x^2-4}{x}}{x+2}$
\end{enumerate}
\end{oefening}

\begin{oefening}[Slim rekenen en schatten]{\oefU\ESS}
Geef eerst zonder rekenmachine een redelijke schatting. Bereken daarna nauwkeurig.
\begin{enumerate}
\item $\dfrac{198\cdot4{,}9}{10{,}2}$
\item $19{,}8^2$
\item $\dfrac{0{,}049\cdot201}{0{,}99}$
\item Een leerling vindt voor $\dfrac{603}{19{,}8}$ de uitkomst $3{,}05$. Leg uit waarom je zonder herberekenen al weet dat dit fout is.
\end{enumerate}
\end{oefening}

% ============================================================
\FVDsection{Machten, wortels en wetenschappelijke notatie}

\begin{oefening}[Machten]{\oefB\ESS}
Vereenvoudig en werk negatieve exponenten weg.
\begin{enumerate}
\item $x^4x^7$
\item $a^3b^2\cdot a^5b$
\item $\dfrac{x^9}{x^4}$
\item $(2x^3)^4$
\item $(3a^{-2}b^3)^2$
\item $\dfrac{12x^5y^{-2}}{4x^2y}$
\item $(2a^2b^{-1})^{-2}$
\item $\dfrac{5x^{-3}}{10x^2}$
\end{enumerate}
\end{oefening}

\begin{oefening}[Wortels]{\oefB\ESS}
Vereenvoudig zo ver mogelijk.
\begin{enumerate}
\item $\sqrt{50}$
\item $\sqrt{108}$
\item $3\sqrt8+2\sqrt{18}$
\item $5\sqrt{12}-\sqrt{75}$
\item $\sqrt6\cdot\sqrt{24}$
\item $(\sqrt7-2)(\sqrt7+2)$
\item $2\sqrt{45}-3\sqrt{20}+\sqrt{80}$
\item $(\sqrt5+\sqrt2)^2$
\end{enumerate}
\end{oefening}

\begin{oefening}[Exact of benaderd?]{\oefU}
Geef telkens aan of een exact antwoord of een decimale benadering het meest geschikt is. Motiveer kort.
\begin{enumerate}
\item De diagonaal van een vierkant met zijde $1$.
\item De lengte van een kabel die je in een winkel moet bestellen.
\item De oppervlakte van een cirkel met straal $3$ in een algebraïsche berekening.
\item Een experimenteel gemeten tijdsduur.
\end{enumerate}
\end{oefening}

\begin{oefening}[Wetenschappelijke notatie]{\oefB\ESS}
Schrijf in wetenschappelijke notatie of als gewoon decimaal getal.
\begin{enumerate}
\item $0{,}000\,0042$
\item $73\,500\,000$
\item $6{,}02\cdot10^{23}$
\item $3{,}4\cdot10^{-7}$
\item $(3\cdot10^5)(4\cdot10^{-8})$
\item $\dfrac{6\cdot10^8}{2\cdot10^{-3}}$
\item $\dfrac{9{,}6\cdot10^{-12}}{3{,}2\cdot10^4}$
\end{enumerate}
\end{oefening}

\begin{oefening}[Ordegrootte in de wetenschappen]{\oefU}
Een bacterie is ongeveer $2\cdot10^{-6}\,\mathrm m$ lang en een menselijke haar ongeveer
$7\cdot10^{-5}\,\mathrm m$ dik.
\begin{enumerate}
\item Hoeveel keer groter is de dikte van de haar dan de lengte van de bacterie?
\item Geef de ordegrootte van beide lengtes.
\item Leg uit waarom wetenschappelijke notatie hier handiger is dan decimale schrijfwijze.
\end{enumerate}
\end{oefening}

% ============================================================
\FVDsection{Procenten, verhoudingen en verandering}

\begin{oefening}[Procentuele verandering]{\oefB\ESS}
Bereken telkens de absolute en procentuele verandering.
\begin{enumerate}
\item $80\to92$
\item $250\to215$
\item $1{,}20\to1{,}38$
\item $640\to800$
\end{enumerate}
\end{oefening}

\begin{oefening}[Verhoudingen]{\oefB}
\begin{enumerate}
\item Een oplossing bevat $18\,\mathrm g$ stof in $300\,\mathrm g$ oplossing. Welk percentage is dat?
\item Op een kaart met schaal $1:25\,000$ is een afstand $7{,}2\,\mathrm{cm}$. Hoe groot is de werkelijke afstand?
\item Een recept voor $4$ personen gebruikt $350\,\mathrm g$ bloem. Hoeveel is nodig voor $7$ personen?
\item Een sensor geeft $2{,}4\,\mathrm V$ bij een bepaalde toestand en $3{,}0\,\mathrm V$ bij een tweede. Met welke factor is de spanning veranderd?
\end{enumerate}
\end{oefening}

\begin{oefening}[Twee keer twintig procent]{\oefU\ESS}
Een toestel kost €500. De prijs stijgt eerst met $20\%$ en daalt daarna met $20\%$.
\begin{enumerate}
\item Bereken de uiteindelijke prijs.
\item Leg uit waarom de prijs niet opnieuw €500 bedraagt.
\item Welke procentuele daling na de stijging zou wel nodig zijn om exact bij €500 uit te komen?
\end{enumerate}
\end{oefening}

% ============================================================
\FVDsection{Algebraïsch rekenen}

\begin{oefening}[Uitwerken]{\oefB\ESS}
Werk uit en vereenvoudig.
\begin{enumerate}
\item $3x(2x-5)$
\item $(x+4)(x-3)$
\item $(2x-3)^2$
\item $(3a+2b)^2$
\item $(5x-2)(5x+2)$
\item $2(x-3)-4(1-2x)$
\end{enumerate}
\end{oefening}

\begin{oefening}[Ontbinden]{\oefB\ESS}
Ontbind zo ver mogelijk.
\begin{enumerate}
\item $5x^2-20x$
\item $x^2-49$
\item $9x^2+30xy+25y^2$
\item $9x^2-36x+36$
\item $18a^3+60a^2b+50ab^2$
\item $x^4-2x^3+x^2$
\item $16x^2-x-2$
\item $2a^4-a^3-a^2$
\end{enumerate}
\end{oefening}

\begin{oefening}[Structuur herkennen]{\oefU\ESS}
Zonder onmiddellijk uit te werken: geef de meest geschikte eerste stap en werk daarna volledig af.
\begin{enumerate}
\item $25x^2-4$
\item $4x^2+12x+9$
\item $6x^3-15x^2+9x$
\item $(x+2)(x^2+1)+(x+2)(x+1)$
\item $a^6+2a^4+a^2$
\end{enumerate}
\end{oefening}

\begin{oefening}[Een fout analyseren]{\oefU}
Een leerling schrijft
\[
(x+3)^2=x^2+9.
\]
\begin{enumerate}
\item Leg precies uit wat ontbreekt.
\item Geef een numerieke controle waarmee je de fout zichtbaar maakt.
\item Geef het correcte resultaat.
\end{enumerate}
\end{oefening}

% ============================================================
\FVDsection{Vergelijkingen, formules en ongelijkheden}

\begin{oefening}[Vergelijkingen]{\oefB\ESS}
Los op en controleer waar zinvol.
\begin{enumerate}
\item $3x-5=16$
\item $4(2x-3)=5x+9$
\item $\dfrac{x-2}{3}+\dfrac{x+1}{2}=5$
\item $(x-3)(x+2)=0$
\item $x^2-9=0$
\item $x^2-5x+6=0$
\end{enumerate}
\end{oefening}

\begin{oefening}[Formules omvormen]{\oefB\ESS}
Maak de gevraagde grootheid vrij.
\begin{enumerate}
\item $v=\dfrac{s}{t}$ naar $s$ en naar $t$.
\item $F=ma$ naar $m$ en naar $a$.
\item $U=RI$ naar $R$ en naar $I$.
\item $A=\pi r^2$ naar $r$.
\item $E_k=\dfrac12mv^2$ naar $v$.
\item $\rho=\dfrac{m}{V}$ naar $V$.
\end{enumerate}
\end{oefening}

\begin{oefening}[Formules in context]{\oefU\ESS}
De periode van een eenvoudige slinger wordt gemodelleerd door
\[
T=2\pi\sqrt{\frac{\ell}{g}}.
\]
\begin{enumerate}
\item Maak $\ell$ vrij.
\item Bereken $\ell$ als $T=2{,}0\,\mathrm s$ en $g=9{,}81\,\mathrm{m/s^2}$.
\item Controleer de eenheid van je resultaat.
\end{enumerate}
\end{oefening}

\begin{oefening}[Eerstegraadsongelijkheden]{\oefB\ESS}
Los op en noteer de oplossing met intervalnotatie.
\begin{enumerate}
\item $2x+3<11$
\item $5-3x\leq14$
\item $4(x-2)>2x+6$
\item $\dfrac{x-1}{3}\geq2$
\item $-4x+7>19$
\item $2(3x-1)\leq4x+10$
\end{enumerate}
\end{oefening}

\begin{oefening}[Tekenschema's]{\oefU\ESS}
Los op met een tekenschema en gebruik de notatie $]\,[$ voor open intervallen.
\begin{enumerate}
\item $(x-2)(x+4)>0$
\item $x^2-9\leq0$
\item $x^2-5x+6\geq0$
\item $-x^2+4x>0$
\item $(2x-3)(x+2)>0$
\item $x(x-2)(x+1)\geq0$
\end{enumerate}
\end{oefening}

\begin{oefening}[Intervallen lezen]{\oefB}
Schrijf als ongelijkheid of als interval.
\begin{enumerate}
\item $x\in]-3,5]$
\item $x\leq4$
\item $-2<x<7$
\item $x\in]-\infty,1[\cup[6,+\infty[$
\item $x\neq3$
\end{enumerate}
\end{oefening}

% ============================================================
\FVDsection{Eenheden en grootheden}

\begin{oefening}[Eenheden omzetten]{\oefB\ESS}
Zet om.
\begin{enumerate}
\item $72\,\mathrm{km/h}$ naar $\mathrm{m/s}$
\item $15\,\mathrm{m/s}$ naar $\mathrm{km/h}$
\item $3{,}5\,\mathrm{cm^2}$ naar $\mathrm{m^2}$
\item $250\,\mathrm{cm^3}$ naar $\mathrm{m^3}$
\item $2{,}4\,\mathrm{g/cm^3}$ naar $\mathrm{kg/m^3}$
\item $5{,}0\,\mathrm{kJ}$ naar $\mathrm J$
\end{enumerate}
\end{oefening}

\begin{oefening}[Dimensie als foutendetector]{\oefU\ESS}
Beoordeel elk resultaat. Leg uit welke fout je op basis van de eenheid of dimensie kunt vermoeden.
\begin{enumerate}
\item Een snelheid wordt gerapporteerd als $18\,\mathrm{m^2/s}$.
\item Een oppervlakte van $35\,\mathrm{cm^2}$ wordt omgezet naar $0{,}35\,\mathrm{m^2}$.
\item Uit $s=vt$ met $v=4\,\mathrm{m/s}$ en $t=5\,\mathrm s$ volgt $s=20\,\mathrm m$.
\item Uit $E=Pt$ met $P$ in watt en $t$ in seconde volgt een energie in joule.
\end{enumerate}
\end{oefening}

% ============================================================
\FVDsection{Rechten en coördinaten}

\begin{oefening}[Vergelijking van een rechte]{\oefB\ESS}
Stel telkens een vergelijking van de rechte op.
\begin{enumerate}
\item rico $4$ en snijpunt met de $y$-as $-1$
\item door $P(2,5)$ met rico $3$
\item door $P(-1,4)$ met rico $-2$
\item door $P(0,3)$ en $Q(2,7)$
\item door $P(1,-2)$ en $Q(4,4)$
\item horizontaal door $P(3,-4)$
\item verticaal door $P(-5,2)$
\end{enumerate}
\end{oefening}

\begin{oefening}[Midden en afstand]{\oefB\ESS}
Bepaal telkens het midden en de afstand tussen de punten.
\begin{enumerate}
\item $A(0,0)$ en $B(3,4)$
\item $A(-2,1)$ en $B(4,9)$
\item $A(5,-3)$ en $B(-1,-3)$
\item $A(-4,-2)$ en $B(2,6)$
\end{enumerate}
Laat wortels exact staan wanneer dat zinvol is.
\end{oefening}

\begin{oefening}[Rechte als model]{\oefU\ESS}
Een experiment levert de meetpunten $(2,7)$ en $(8,25)$ op.
\begin{enumerate}
\item Stel de vergelijking op van de rechte door beide punten.
\item Interpreteer de richtingscoëfficiënt als verandering van $y$ per eenheid $x$.
\item Voorspel $y$ voor $x=10$.
\item Waarom moet je voorzichtig zijn wanneer je dezelfde rechte gebruikt voor $x=1000$?
\end{enumerate}
\end{oefening}

% ============================================================
\FVDsection{Pythagoras en basisgoniometrie}

\begin{oefening}[Pythagoras]{\oefB\ESS}
Bereken de ontbrekende zijde.
\begin{enumerate}
\item Rechthoekszijden $6$ en $8$.
\item Schuine zijde $13$ en één rechthoekszijde $5$.
\item Rechthoekszijden $7$ en $11$.
\item Een rechthoek heeft lengte $12$ en breedte $5$. Bereken de diagonaal.
\end{enumerate}
\end{oefening}

\begin{oefening}[Sinus, cosinus en tangens]{\oefB\ESS}
\begin{enumerate}
\item Schuine zijde $10$, aanliggende zijde bij $\alpha$ gelijk aan $6$: bepaal $\cos\alpha$.
\item Overstaande zijde $9$, aanliggende zijde $12$: bepaal $\tan\alpha$.
\item $\sin\alpha=\frac35$ en de schuine zijde is $20$: bepaal de overstaande zijde.
\item $\tan\alpha=\frac43$ en de aanliggende zijde is $15$: bepaal de overstaande zijde.
\item Een rechthoekige driehoek heeft rechthoekszijden $8$ en $15$. Bepaal de kleinste scherpe hoek.
\end{enumerate}
\end{oefening}

\begin{oefening}[STEM: hellingshoek]{\oefU\ESS}
Een helling stijgt $1{,}8\,\mathrm m$ over een horizontale afstand van $12\,\mathrm m$.
\begin{enumerate}
\item Bereken de lengte van de helling.
\item Bereken de hellingshoek met de horizontale.
\item De helling wordt beschreven als ``$15\%$''. Onderzoek hoe dat percentage verband houdt met je driehoek.
\end{enumerate}
\end{oefening}

\begin{oefening}[Afstand zonder formule uit het hoofd]{\oefU}
De punten zijn $A(-3,4)$ en $B(5,-2)$.
\begin{enumerate}
\item Teken schematisch de horizontale en verticale verschillen.
\item Leid met Pythagoras de afstand $|AB|$ af.
\item Vergelijk je afleiding met de algemene afstandsformule.
\end{enumerate}
\end{oefening}

% ============================================================
\FVDsection{Welk gereedschap heb je nodig?}

\begin{oefening}[Kies eerst, reken daarna]{\oefU\ESS}
Schrijf bij elk probleem eerst op welk gereedschap je kiest. Los daarna op.
\begin{enumerate}
\item Een zwemclub huurt een bus. De vaste kost is €180 en per leerling komt €12 bij. Stel $K(x)$ op en bepaal voor hoeveel leerlingen de kost kleiner is dan €600.
\item Een rechthoekige tuin heeft breedte $x$ meter en lengte $2x+3$ meter. Stel $O(x)$ op. Voor welke $x$ is de omtrek $66$ meter?
\item Een taxi rekent €4 opstapgeld en €2,50 per kilometer. Voor welke afstanden is de prijs hoogstens €31,50?
\item Een kostfunctie verloopt lineair. Na $10$ minuten is de kost €140 en na $25$ minuten €320. Stel het voorschrift op en bereken de kost na $21$ minuten.
\end{enumerate}
\end{oefening}

\begin{oefening}[Combineren]{\oefU\ESS}
Een drone vliegt in een rechte lijn van $A(0,0)$ naar $B(120,50)$, met coördinaten in meter.
De vlucht duurt $10\,\mathrm s$.
\begin{enumerate}
\item Bereken de afgelegde rechte afstand.
\item Bereken de gemiddelde grootte van de snelheid in $\mathrm{m/s}$.
\item Zet die snelheid om naar $\mathrm{km/h}$.
\item Een tweede vlucht verloopt gemiddeld aan $15\,\mathrm{m/s}$. Bereken de procentuele toename.
\item Noteer bij elke stap welk wiskundig gereedschap je gebruikte.
\end{enumerate}
\end{oefening}

\begin{oefening}[Van formule naar betekenis]{\oefU}
De kinetische energie is
\[
E_k=\frac12mv^2.
\]
Voor een voorwerp met massa $0{,}80\,\mathrm{kg}$ wordt $E_k=90\,\mathrm J$ gemeten.
\begin{enumerate}
\item Bereken de snelheid.
\item Een leerling vindt $v=225\,\mathrm{m/s}$. Geef twee manieren om te controleren dat dit niet klopt.
\item Wat gebeurt er met $E_k$ als de snelheid verdubbelt? Leg uit zonder nieuwe getallen te kiezen.
\end{enumerate}
\end{oefening}

% ============================================================
\FVDsection{Gevorderde uitdagingen}

\begin{oefening}[Een parameterprobleem]{\oefG}
Voor
\[
f_a(x)=ax+4
\]
gaat de grafiek door het punt $P(3,10)$.
\begin{enumerate}
\item Bepaal $a$.
\item Bepaal het snijpunt met de $x$-as.
\item Voor welke $x$ geldt $f_a(x)>0$?
\item Interpreteer je antwoord grafisch.
\end{enumerate}
\end{oefening}

\begin{oefening}[Ontwerpen met voorwaarden]{\oefG}
Een rechthoekig zonnepaneel heeft oppervlakte $1{,}80\,\mathrm{m^2}$.
De lengte is $0{,}60\,\mathrm m$ groter dan de breedte.
\begin{enumerate}
\item Stel een vergelijking op voor de breedte.
\item Bepaal de afmetingen.
\item Bereken de diagonaal.
\item Een transportkist laat maximaal een diagonaal van $1{,}75\,\mathrm m$ toe. Past het paneel? Motiveer.
\end{enumerate}
\end{oefening}

\begin{oefening}[Model kritisch beoordelen]{\oefG}
Een sensor wordt tussen twee kalibratiepunten lineair gemodelleerd:
\[
(0,0{,}20)\qquad\text{en}\qquad(100,4{,}80).
\]
\begin{enumerate}
\item Stel het lineaire model op.
\item Bereken de voorspelde sensorwaarde bij invoer $60$.
\item Bereken welke invoer hoort bij sensorwaarde $3{,}00$.
\item Het model voorspelt ook een waarde bij invoer $1000$. Bereken die waarde, maar bespreek vervolgens waarom de berekening niet automatisch een betrouwbare fysische voorspelling is.
\end{enumerate}
\end{oefening}


\end{document}
